\chapter*{Tóm tắt}
\label{summary}
\begin{spacing}{1.15}
Phân đoạn ảnh y sinh đóng vai trò then chốt trong hỗ trợ chẩn đoán và điều trị bệnh, đặc biệt là trong việc nhận diện vùng tổn thương da (dermoscopy) hay cấu trúc tuyến trong ảnh mô học (histology). Mặc dù kiến trúc U-Net đã trở thành tiêu chuẩn cho bài toán này, khả năng biểu diễn của nó vẫn bị giới hạn bởi nền tảng tính toán cổ điển, đặc biệt khi xử lý các tập dữ liệu nhỏ dễ gây quá khớp. Khóa luận đề xuất HQMoSS-Net -- một kiến trúc mạng nơ-ron lai lượng tử -- cổ điển trên nền U-Net. Điểm trọng tâm của HQMoSS-Net là việc tích hợp tại vùng nút thắt cổ chai một khối Quantum Mixture of Experts (QuantumMoE) đa tỷ lệ -- trong đó mỗi nhánh xử lý là một mạch lượng tử tham số hóa dùng mã hóa biên độ và xử lý mảnh ảnh không gian ở một tỷ lệ khác nhau -- kết hợp với một khối chuỗi quét chọn lọc lấy cảm hứng từ Mamba (Lightweight Selective Scan) để tổng hợp phụ thuộc tầm xa với chi phí tuyến tính, và một khối tích chập đa tỷ lệ (Residual Dilated Multi-Scale) ở bộ khung cổ điển. Mô hình được huấn luyện đầu-cuối bằng hàm mất mát Binary Cross-Entropy trong môi trường giả lập lượng tử (PennyLane), và đánh giá bằng kiểm định chéo $5$-fold trên hai tập dữ liệu PH2 và GlaS 2015, so sánh với năm mô hình cổ điển (U-Net, UNet++, Attention U-Net, V-Net, R2U-Net). Kết quả cho thấy HQMoSS-Net đạt Soft IoU $0.8497$ trên PH2 và $0.7038$ trên GlaS 2015, vượt trội mọi mô hình cơ sở dù có số tham số nhỏ hơn nhiều lần; phân tích bóc tách khẳng định đóng góp bổ trợ của từng thành phần, một đối chứng với khối Classical MoE cùng kiến trúc nhưng nhiều tham số hơn cho tín hiệu khiêm tốn rằng lợi ích không chỉ đến từ cấu trúc, và bước trực quan hóa Seg-Grad-CAM minh họa mô hình tập trung đúng vào vùng tổn thương.
\end{spacing}